\section{Space, reachability, complementation, and finite games}
\label{sec:space}

This section gives explicit finite forms of the space-complexity theorems.
The underlying classical results are due to Savitch and, for
complementation, independently to Immerman and Szelepcs\'enyi
\cite{Savitch1970,Immerman1988,Szelepcsenyi1988}.
Their interpretation here is more specific than an appeal to those results:
we give bounded algorithms, witness transformations, and the Boolean
invariants that permit quantifier-free induction.  The presentation of
space-bounded computation and quantified Boolean formulae is also consistent
with the author-hosted text of Arora and Barak \cite{AroraBarak}.

\subsection{Finite graph conventions and path elimination}

A graph has vertices $0,\ldots,N-1$, where $N\geq1$, and an adjacency bit
$E(u,v)$.  An explicit graph stores these $N^2$ bits; an implicit graph
provides a fixed bounded algorithm for the bit.  Throughout, $a,b<N$.
A walk of length $j$ is a list $v_0,\ldots,v_j<N$ such that
$E(v_i,v_{i+1})=1$ for every $i<j$.  Empty walks are allowed.  A list of
$j+1$ vertices can be encoded by $j+1$ base-$\max(2,N)$ digits, with a
separate length field.  Consequently every path-existence assertion with
$j\leq B$ is bounded by an explicitly primitive-recursive number, for
example $(B+2)\max(2,N)^{B+1}$ after a routine tuple encoding.

For a fixed source $a$, define reachability bits by
\[
 r_0(v)=[v=a],\qquad
 r_{i+1}(v)=r_i(v)\vee
       \bigvee_{u<N}\bigl(r_i(u)\wedge E(u,v)\bigr).
 \tag{S.1}
\]
The entire row $(r_i(0),\ldots,r_i(N-1))$ is an $N$-bit integer.
Equation (S.1), evaluated by nested bounded folds, is a primitive recursion
on $i$.  This definition gives a total numerical function even when a
later space-efficient implementation computes its values by a different
algorithm.

\begin{theorem}[S01: bounded reachability and constructive path shortening]
\label{thm:S01}
The bit $r_i(v)$ is one exactly when a walk from $a$ to $v$ of length at
most $i$ exists.  A witnessing walk can be obtained by primitive recursion
and bounded search.  Every walk in an $N$-vertex graph has a walk with the
same endpoints and length at most $N-1$.  Thus $r_{N-1}=r_N$, and their
common value decides unrestricted finite-walk reachability.
\end{theorem}
\begin{proof}
At $i=0$ the sole possible walk is the empty walk at $a$.  If (S.1) is
one at stage $i+1$, either the old bit is one, in which case retain its
walk, or select the least $u<N$ with $r_i(u)E(u,v)=1$ and append $v$ to
the already constructed walk to $u$.  Conversely, a nonempty walk of
length at most $i+1$ has a last predecessor reached in at most $i$
steps; a walk of length zero is retained by the first disjunct.  This
proves both directions with actual witnesses.

For shortening, scan a supplied walk from left to right while maintaining
a list of distinct vertices with the same current endpoint.  When the
next vertex is new, append it.  When it already occurs at position $j$,
delete the suffix after that occurrence.  The surviving prefix is still a
walk: in the repeated-vertex case its last vertex is precisely the vertex
just reached, so the next input edge, if any, remains usable.  A distinct
list of vertices below $N$ has length at most $N$.  One finite proof is to
form its $N$ incidence bits: when a new distinct vertex is inserted, their
sum increases by one, and each incidence bit is at most one.  Therefore
the final distinct walk has at most $N-1$ edges.  Every operation is a
bounded list scan or update.  Shortening a walk of length at most $N$
proves $r_N\leq r_{N-1}$; the reverse inclusion follows directly from
(S.1).  Shortening arbitrary supplied walks proves the final assertion.

For the arithmetic proof, fix an outer bound $B$ on the supplied walk
length.  The invariant says, for every prefix length at most $B$, that
the current list is distinct, consists of vertices below $N$, has the
same endpoint as the input prefix, and has valid consecutive edges.  All
its quantifiers are bounded by $B,N$ and the displayed list-code bound.
Its Boolean compilation is a legitimate quantifier-free induction
formula.  The path-construction invariant is handled in the same way.
\end{proof}

\subsection{Configuration graphs and logarithmic space}

Fix a machine with $q$ control states, $k$ one-sided work tapes, and
work alphabet of size $g$.  Its input tape is read-only.  A space budget
$s\geq1$ permits the first $s$ cells of each work tape; crossing the
budget is made rejecting.  A configuration records the control state,
the input-head position, the $k$ work-head positions and the work-tape
contents.  The number of possible ordinary configurations is at most
\[
 Q=q(n+2)(s+1)^k g^{ks}.                 \tag{S.2}
\]
The harmless $+1$ in each head factor permits a boundary sentinel.
Add a distinguished accepting sink and rejecting sink, and let $\ell$
be an explicitly computed bit length sufficient for every such record.
The vertex set can be padded to $N=2^\ell$: invalid records have no
outgoing edges.  A transition from any accepting configuration to the
accepting sink supplies a unique target.  For this fixed machine,
\[
 \ell\leq C\bigl(s+\lceil\log_2(n+2)\rceil+1\bigr)
 \tag{S.3}
\]
for a fixed natural constant $C$.  We use integer bit length; no real
logarithm is needed in the arithmetic definitions.

\begin{theorem}[S02: the finite configuration theorem]
\label{thm:S02}
Acceptance within the supplied space budget is equivalent to reachability
of the accepting sink in the above graph, and hence to existence of a
computation of fewer than $N$ transitions in that graph.  The
configuration graph and its edge test are primitive-recursive.  For
$s\geq\lceil\log_2(n+2)\rceil$, its vertices and edge test use $O(s)$
working bits.  A deterministic computation which eventually halts has
halted before repeating an ordinary configuration.
\end{theorem}
\begin{proof}
Write down the successive configurations of any accepting computation;
this is a graph walk, to which S01 applies.  Conversely, a graph walk
from the initial record to the accepting sink lists legal transitions,
so deleting the final sink edge gives an accepting computation.  Validity
of each record and each transition is decided by scanning the finite
transition table, checking head movements, and comparing unchanged tape
cells.  A scan needs only the two configuration records and bounded
indices.  Equation (S.3) therefore proves the claimed space bound.

If a deterministic computation repeats a configuration, its future is
repeated as well: induction on the number of further transitions proves
equality of the corresponding configurations.  A first halting time
cannot lie strictly after such a repetition.  This is a finite statement
about a supplied halting trace; it does not invoke a decision procedure
for unbounded halting.  Alternatively one can normalize any deterministic
space-bounded decider by rejecting at the explicit configuration clock.
The proof of equivalence uses its stated totality assumption.
\end{proof}

Whenever a complexity bound $s(n)$ is used operationally below, it is a
fixed primitive-recursive bound together with an $O(s(n))$-space procedure
for obtaining or marking the budget.  Equivalently, the theorem may be
read with an explicit budget parameter.  This qualification avoids silently
assuming that an arbitrary externally described numerical bound can be
computed in the advertised space.  The standard logarithmic, polynomial,
and exponential space bounds have the required constructors.

\begin{theorem}[S03: directed reachability is NL-complete]
\label{thm:S03}
Directed graph reachability is complete for $\mathrm{NL}$ under
logarithmic-space many-one reductions.  In particular,
\[
 \mathrm L\subseteq\mathrm{NL}\subseteq\mathrm P.
\]
\end{theorem}
\begin{proof}
For membership, nondeterministically choose a walk of at most $N-1$
edges, retaining only the current vertex, target, and step counter.
An explicit input graph can be rescanned to test an edge.  These
registers use $O(\log(N+2))$ bits.  S01 proves both acceptance directions.

For hardness, fix a logarithmic-space machine recognizing a language.
In (S.3), $\ell=O(\log(n+2))$, hence $N=2^\ell$ is polynomial in $n$.
Output its configuration graph in a fixed adjacency-matrix encoding.
The transducer runs through the pair of vertex codes $(u,v)$, computes
the corresponding edge bit, and writes it.  It stores only two
$O(\log n)$-bit records and bounded scan indices; it never stores the
whole graph.  The output length $N^2$ is polynomial.  S02 proves the
reduction equation.  Start and accepting vertices are included in the
output in the same space.

Deterministic logarithmic-space computation is a special case of
nondeterministic computation.  For the second containment, an explicit
graph can be solved deterministically by the $N$ iterations of (S.1),
using $O(N^3)$ elementary adjacency operations.  For the polynomial-size
configuration graph just constructed, the resulting time, including
edge generation or explicit graph construction, is polynomial in $n$.
The transducers and their simulation equations provide numerical
compiler witnesses for the two containments.
\end{proof}

\begin{corollary}[S04: even acyclic directed reachability is NL-complete]
\label{thm:S04}
The restriction of directed reachability to directed acyclic graphs is
NL-complete, by a logarithmic-space reduction whose correctness has
explicit path maps.
\end{corollary}
\begin{proof}
Given $G$ with $N$ vertices, form vertices $(v,i)$ for $v<N$ and
$0\leq i<N$.  Add $(u,i)\to(v,i+1)$ exactly when $i<N-1$ and either
$u=v$ or $E(u,v)=1$.  Every edge increases the second coordinate, so
no directed cycle is possible.  A path from $a$ to $b$ of length at
most $N-1$ can be padded by stays to give a path from $(a,0)$ to
$(b,N-1)$.  Conversely erase stay edges in a layered path.  For $N=1$
this is just the length-zero test.  The two bounded counters and a
rescan of $G$ give a logarithmic-space transducer.  Membership follows
from S03, because the general reachability algorithm applies.
\end{proof}

\subsection{Savitch's recursion and its actual workspace}

For $a,b<N$ define Boolean matrices $R_j$ by
\[
 \begin{split}
 R_0(a,b)&=[a=b]\vee E(a,b),\\
 R_{j+1}(a,b)&=\bigvee_{z<N}
                 \bigl(R_j(a,z)\wedge R_j(z,b)\bigr).
 \end{split}                                             \tag{S.4}
\]
As a numerical definition this is primitive recursion on the whole
$N^2$-bit matrix.  This observation establishes primitive-recursive
status; storing that matrix is \emph{not} the space-efficient algorithm.

\begin{theorem}[S05: Savitch's finite theorem]
\label{thm:S05}
The bit $R_j(a,b)$ is one exactly when there is a walk from $a$ to $b$
of length at most $2^j$.  There is a deterministic implementation of
$R_j(a,b)$ using
\[
 O\bigl((j+1)(\lceil\log_2(N+2)\rceil+
                    \lceil\log_2(j+2)\rceil)+w\bigr)
 \tag{S.5}
\]
working bits when the edge test uses $w$ working bits.  Its time is
bounded by an explicit primitive-recursive function; for a polynomial-time
edge test one may use
$C(2N+2)^{j+2}(\ell+j+2)^C$ with a suitable fixed $C$.
\end{theorem}
\begin{proof}
The assertion at $j=0$ is the definition of a walk of length at most
one.  Concatenate the two walks certified by a successful midpoint in
(S.4) to prove soundness at $j+1$.  Conversely split a walk after
$\min(2^j,\text{its length})$ edges.  Both pieces have length at most
$2^j$, so their common midpoint makes (S.4) true.  These operations
are bounded list transformations.  If a witness is wanted, search for
the least successful midpoint and recursively concatenate the witnesses.

For the space claim, implement the recurrence by depth-first evaluation.
A frame stores its endpoints, a midpoint counter, its level, and a
constant-size phase indicating whether it is trying a new midpoint or
returning from one of the two child calls.  If the first child is false,
advance the midpoint.  If it is true, evaluate the second child.  Return
true on the first pair of successful children, and false after the last
midpoint.  A parent need not store either child's recursion stack after
that child returns.  There are at most $j+1$ active frames.  Each has
$O(\log(N+2)+\log(j+2))$ bits; edge-test scratch space can be reused
at the leaves.  This proves (S.5).

Every frame has at most $2N$ child calls.  Bounding all control and
edge-test costs by the displayed polynomial gives the stated coarse
time bound by induction.  For arbitrary specified bounded edge code,
replace that polynomial by its supplied primitive-recursive clock.
The explicit time bound also allows the stack machine to be defined by
clocked primitive recursion, without using general well-founded recursion
as an unexamined arithmetic principle.

For the iPRA induction, fix an outer maximum $J$.  At levels $j\leq J$,
all walk witnesses are bounded by a base-$\max(2,N)$ code of length
$2^J+1$.  Thus the invariant universally quantifying over vertices,
levels and witnesses is bounded.  The stack-correctness invariant is
bounded by the displayed time and stack-code bounds.  Both compile to
Boolean equations for quantifier-free induction.
\end{proof}

\begin{corollary}[S06: the principal space equalities]
\label{thm:S06}
For the explicit budget conventions above and $s(n)\geq\log_2(n+2)$,
\[
 \mathrm{NSPACE}(s)\subseteq\mathrm{DSPACE}(s^2).
\]
Consequently,
\[
 \begin{gathered}
 \mathrm{NL}\subseteq\mathrm{DSPACE}(\log^2 n),\qquad
 \mathrm{NPSPACE}=\mathrm{PSPACE},\\
 \mathrm{NEXPSPACE}=\mathrm{EXPSPACE},\qquad
 \mathrm{PSPACE}\subseteq\mathrm{EXPTIME}.
 \end{gathered}
\]
\end{corollary}
\begin{proof}
Use the configuration graph with $N=2^\ell$ and evaluate
$R_\ell(a,b)$.  S01 says that this reaches every possible accepting
configuration, while (S.3) and (S.5) give $O(s^2)$ working bits.
The logarithmic corollary follows by substitution.  The square of a
polynomial is a polynomial.  The square of $2^{p(n)}$ is $2^{2p(n)}$,
which is still an exponential-space bound.  The reverse inclusions in
the two equalities hold because deterministic machines are special
nondeterministic machines.  Finally a deterministic polynomial-space
decider has only $2^{\poly(n)}$ configurations and cannot repeat one
before halting, by S02; simulation up to that clock is exponential time.
The same normalization also yields
$\mathrm{DSPACE}(s)\subseteq\mathrm{DTIME}(2^{O(s)})$ for the stated
fixed machines and budget constructors.
\end{proof}

\subsection{Inductive counting without negative guesses}

Write $A_i=\{v<N:r_i(v)=1\}$ and $c_i=\sum_{v<N}r_i(v)$.  These are
notation for a bit table and its bounded sum, not set-existence axioms.
A guessed path of length at most $i$ is checked online: retain its
current vertex and counter, guess the next vertex, check its edge, and
at a guessed stopping time require the desired endpoint.  Every branch
is stopped by the bound $i$.

\begin{lemma}[S07: forced enumeration]
\label{thm:S07}
Suppose the exact number $c_i$ is supplied.  There is a nondeterministic
$O(\log(N+2))$-space procedure which, on a requested $v$, certifies
$v\notin A_{i+1}$.  Every completed successful branch certifies that
fact, and if the fact holds a successful branch can be constructed by a
primitive-recursive procedure.  Here $0\leq i<N$.
\end{lemma}
\begin{proof}
Scan $u=0,\ldots,N-1$ in order, keeping a counter $d=0$.
At each $u$ guess whether to select it.  A selected vertex must be
accompanied by a checked walk from $a$ of length at most $i$; also check
$u\ne v$ and $E(u,v)=0$, then increase $d$.  An unselected vertex
requires no claim.  At the end require $d=c_i$.

Every selected vertex belongs to $A_i$, and no vertex is selected twice.
A subset of the finite $0/1$ table $r_i$ with the same sum $c_i$ equals
that table: if an omitted table entry were one, summing the componentwise
inequalities would give $d<c_i$.  Hence a successful branch has checked
$u\ne v$ and $E(u,v)=0$ for every $u\in A_i$.  Formula (S.1) now says
$v\notin A_{i+1}$.

Conversely, if $v\notin A_{i+1}$, select exactly the vertices in
$A_i$, and for each select the canonical path constructed in S01.
All checks pass, and exactly $c_i$ vertices are selected.  Computing this
successful transcript may use large temporary tables; the recognizing
machine only guesses it and checks paths one at a time.  This distinction
is essential.  The recognizing machine stores $i,v,u,d,c_i$ and a few
path registers, all logarithmic in $N$; it does not store $A_i$.
\end{proof}

\begin{theorem}[S08: constructive space complementation]
\label{thm:S08}
Nonreachability has a nondeterministic logarithmic-space recognizer whose
accepting transcripts are sound and whose successful transcript can be
constructed primitively recursively whenever nonreachability holds.
For the budgeted machine presentations above,
\[
 \mathrm{NSPACE}(s)=\mathrm{coNSPACE}(s)
 \quad(s\geq\log_2(n+2)),
 \qquad \mathrm{NL}=\mathrm{coNL}.
\]
These classes are closed under finite unions, intersections, differences,
and symmetric differences.
\end{theorem}
\begin{proof}
Start from $c_0=1$.  Given the correct $c_i$, compute a candidate
$c_{i+1}$ by scanning every $v<N$.  Guess either ``reachable'' or
``not reachable.''  In the first case verify a path of length at most
$i+1$ and increment the new counter.  In the second case run S07.
All unsuccessful guesses reject the branch.  Every branch completing
the scan has classified each vertex correctly, so its new count equals
$c_{i+1}$.  At least one branch completes: use the actual finite bits
$r_{i+1}(v)$, canonical positive paths, and the successful transcripts
from S07.  Induction on $i=0,\ldots,N-1$ yields the exact $c_N$ on
every branch which survives all scans.

For the final target $b$, enumerate vertices once more, selecting only
vertices accompanied by a path of length at most $N$, requiring every
selected vertex to differ from $b$, and requiring exactly $c_N$
selections.  The equal-cardinality argument forces enumeration of all
of $A_N$.  Therefore the final test can succeed exactly when
$r_N(b)=0$, which by S01 is nonreachability.  It never accepts a bare
unsupported assertion of nonreachability.

There are at most $N$ stages, $N$ queried vertices per stage, $N$
forced-enumeration vertices per negative query, and $N$ edges per checked
path.  Thus $O(N^4)$ edge tests suffice on every branch, apart from
polynomial control overhead.  Only a constant number of indices, counts,
and current vertices are live at once.  This proves logarithmic space
for explicit graphs.  For the configuration graph, every vertex has
$\ell=O(s)$ bits and its edge test has $O(s)$ workspace; replacing an
explicit adjacency lookup by that test proves the general statement.
The accepting transcript bound is exponential in a primitive-recursive
bound and remains primitive recursive.

The iPRA induction invariant is stronger than ``the count is correct''
on one favored computation.  It states that every transcript prefix
within the explicit branch-code bound which completes stage $i$ has
count $\sum_{v<N}r_i(v)$, and that the canonical transcript constructor
completes that stage.  All transcript, vertex, path and stage quantifiers
are bounded.  Its Boolean compilation justifies the induction and
produces positive witnesses for completeness.

For closure under union, first choose which machine to run.  For
intersection, run the two normalized recognizers sequentially and
accept only if both accept, reusing their workspace.  Given accepting
traces for both, concatenation gives an accepting combined trace;
projection gives the converse.  Complementation, intersection and union
then construct differences and symmetric differences.  Reapplying the
complement compiler gives the reverse class inclusion, so the displayed
equality is constructive at the level of presentation compilers.
\end{proof}

\begin{corollary}[S09: 2SAT is NL-complete]
\label{thm:S09}
Satisfiability of formulae with clauses of size at most two is
NL-complete under logarithmic-space reductions.  Its restriction to
clauses of the forms $\neg x_u\vee x_v$ and unit clauses is already
NL-hard.
\end{corollary}
\begin{proof}
The constructive criterion in Theorem~\ref{thm:N13} says that a 2CNF
formula is unsatisfiable precisely when some literal and its negation
are mutually reachable in its implication graph.  Guess the variable
and the two bounded paths.  This is an NL recognizer for 2CNF
unsatisfiability, since the implication graph's edge bits can be
computed by rescanning the formula.  S08 yields an NL recognizer for
satisfiability.

For hardness, transform a graph with source $a$ and target $b$ into
\[
 F_G=x_a\ \wedge\ \neg x_b\ \wedge\!
              \bigwedge_{E(u,v)=1}(\neg x_u\vee x_v).
 \tag{S.6}
\]
If $a$ reaches $b$, the implication clauses propagate truth from $x_a$
to $x_b$, contradicting the unit clause $\neg x_b$.  If $a$ does not
reach $b$, the assignment $x_v=r_N(v)$ satisfies every clause: a true
reachable tail of an edge has a reachable head, $x_a=1$, and $x_b=0$.
Thus nonreachability reduces to 2SAT with explicit witness extraction.
The formula can be streamed with two vertex counters.  S03 and S08
make nonreachability NL-hard: reduce the complement of an arbitrary NL
language to reachability and then complement the resulting equality.
All complements here are complements of specified Boolean predicates.
This proves the result without an unbounded excluded-middle step.
\end{proof}

\subsection{Quantified Boolean formulae with controlled syntax size}

A Boolean formula uses variables, negation, conjunction, disjunction,
and the finite quantifiers $\exists z\in\{0,1\}$ and
$\forall z\in\{0,1\}$.  Evaluation of these quantifiers means respectively
Boolean disjunction and conjunction of the two evaluations.  A closed
formula has a determinate evaluation bit by a bounded recursive traversal
of its finite syntax.  Malformed syntax is assigned value zero and
excluded by a validity bit.  The language TQBF consists of valid closed
formulae with value one.

\begin{theorem}[S10: bounded evaluation of QBF]
\label{thm:S10}
There is a primitive-recursive evaluator for closed Boolean formulae,
and a deterministic polynomial-space implementation.  Truth has an
explicit finite winning strategy for existential choices; falsity has
an explicit finite counterstrategy for universal choices.
\end{theorem}
\begin{proof}
Use depth-first syntax evaluation.  Store the current assignment,
a syntax-position stack, and a phase bit at each quantifier.  At an
existential quantifier evaluate both assignments in order until one
is true; at a universal quantifier require both to be true.  After a
child returns, restore the old value of the bound variable.  For syntax
length $m$, there are at most $m$ stack frames; positions and variable
indices use $O(\log(m+2))$ bits, and an assignment uses at most $m$
bits.  Even a deliberately wasteful implementation copying the assignment
at each frame uses $O(m^2+m\log(m+2))$ bits.  A coarse clock
$2^{Cm}(m+2)^C$ suffices: a branch has depth at most $m$, every node
has at most two children, and each local operation is polynomial.
Running this implementation to that clock is a fixed primitive recursion.

For a true existential node choose the least child with value one;
for a true universal node retain both true children.  At a false
universal node choose the least false child, and at a false existential
node retain both false children.  The recursive construction is bounded
by the evaluation tree.  Induction on its finite depth proves that the
resulting strategy wins against every opposing sequence of choices.
This supplies witnesses and counterwitnesses without invoking an
infinite game-determinacy principle.
\end{proof}

The next construction is needed because writing down both halves of
(S.4) recursively would produce an exponentially long formula.
Let $U,V,M,A,B$ denote vectors of $\ell$ Boolean variables, and let
$D_x(U,V)$ be a Boolean formula for equality or one configuration step
of the fixed machine on input $x$.  Such a formula has polynomial size
in $n+s$.  Explicitly, enumerate its finitely many transition rules
and all possible head positions, require the stated control-state and
scanned-symbol bits, prescribe the changed head and symbol fields,
and require equality of every untouched cell.  The number of work
heads and the transition table are fixed, so this enumeration is
polynomial.  Input symbols are the fixed bits of $x$.  Validity checks
exclude malformed records; sink transitions are added explicitly.

Define formulae with free endpoint vectors by
\[
 \begin{split}
 \Theta_0(U,V)&=D_x(U,V),\\
 \Theta_{j+1}(U,V)
 &=\exists M\ \forall z\ \exists A\ \exists B\,
       \bigl(H_z(U,V,M,A,B)\wedge\Theta_j(A,B)\bigr),
 \end{split}                                               \tag{S.7}
\]
where every displayed vector quantifier abbreviates $\ell$ Boolean
quantifiers, all newly bound variables are fresh, and
\[
 H_z=\bigl(\neg z\to(A=U\wedge B=M)\bigr)\wedge
        \bigl(z\to(A=M\wedge B=V)\bigr).
 \tag{S.8}
\]
Vector equality is a conjunction of bit equalities; implication in
(S.8) is the ordinary Boolean connective $\neg P\vee Q$.

\begin{theorem}[S11: polynomial-size quantified reachability]
\label{thm:S11}
For all endpoint assignments, $\Theta_j(U,V)$ has value $R_j(U,V)$.
Its syntax has one copy of $D_x$ and only $O(j\ell)$ extra connective
and variable occurrences.  With binary variable indices its encoded
length is bounded by
\[
 O\bigl((|D_x|+j\ell)\log((j+1)\ell+2)\bigr),
 \tag{S.9}
\]
where $|D_x|$ counts syntax nodes before variable-name encoding.  Construction,
renaming, and substitution of constant endpoints take polynomial time
in the output bound.
\end{theorem}
\begin{proof}
Fix $M$.  At $z=0$, the constraints in (S.8) force $A=U,B=M$;
at $z=1$ they force $A=M,B=V$.  These assignments exist explicitly
and are unique.  Thus the two universal branches are true exactly
when both $\Theta_j(U,M)$ and $\Theta_j(M,V)$ are true.
The outer existential quantifier then takes the disjunction over
$M<N$, exactly as in (S.4).  This proves the recurrence by Boolean
calculation, and induction proves the evaluation identity.  In
particular the midpoint must be chosen \emph{before} the universal
selector; choosing it separately after the selector would invalidate
the proof.

There is only one textual occurrence of $\Theta_j$ on the right of
(S.7).  Each level adds three vector blocks, one selector, and a
constant number of vector equalities.  These have $O(\ell)$ nodes.
The total number of fresh variables is $O((j+1)\ell)$, so variable names
cost $O(\log((j+1)\ell+2))$ bits.  This factor applies to the base
formula too, because its endpoint variables are renamed at the deepest
level.  This proves (S.9).
A bounded traversal renames free endpoint variables without expanding
or copying any recursive subformula.  All loops have polynomial bounds.
If a prenex convention is desired, pull the fresh inner quantifiers
outward across the conjunctions in (S.7); their variables are absent
from the corresponding guard.  Both directions of each transformation
are verified by the two Boolean quantifier cases.  This introduces no
duplication and no superpolynomial growth.
\end{proof}

\begin{theorem}[S12: TQBF is PSPACE-complete]
\label{thm:S12}
Every fixed polynomial-space machine presentation has a polynomial-time
many-one reduction to TQBF, with an iPRA proof of its numerical
correctness equation.  Together with S10 this proves PSPACE completeness.
\end{theorem}
\begin{proof}
For the polynomial-space machine and input $x$, compute its polynomial
space bound, configuration length $\ell$, initial record $a$, accepting
sink $b$, and the polynomial-size formula $D_x$.  Output the closed
formula $\Theta_\ell(a,b)$.  S11 says its value is $R_\ell(a,b)$;
S01 and S02 say that this is exactly the input's acceptance bit.
Here $\ell=O(\poly(n))$, so (S.9) is polynomial in $n$ and the
construction is polynomial time.  Every compiler loop, equality test,
and syntax substitution has an explicit polynomial bound.  Numerical
correctness is obtained by the bounded invariants already proved;
this is the stated non-formalized iPRA proof, not a claim that a Lean
object-derivation has already been generated.  S10 proves membership.
\end{proof}

\subsection{Succinct games with a polynomial horizon}

A useful immediate consequence is a precise game-complexity theorem.
An instance supplies a state width $r$, Boolean circuits
$\operatorname{owner}(q)$, $\operatorname{next}(q,b)$ and
$\operatorname{win}(q)$, an initial $r$-bit state, and a horizon $h$
in unary.  Players choose one bit per move.  The owner circuit says
whether the existential or universal player chooses, the next circuit
updates the state, and after exactly $h$ moves the existential player
wins iff $\operatorname{win}(q)=1$.  Earlier terminal positions can be
made absorbing.  Circuits are explicit finite acyclic circuits; their
outputs can be evaluated with polynomial workspace.

\begin{theorem}[S13: finite strategies and game completeness]
\label{thm:S13}
The winner of a supplied game is decided by a primitive-recursive
Boolean function.  One of the players has a canonical finite winning
strategy.  Deciding whether the existential player wins is PSPACE-complete
under polynomial-time reductions.
\end{theorem}
\begin{proof}
Write $W_0(q)=\operatorname{win}(q)$.  Given $W_i$, put
\[
 W_{i+1}(q)=
 \begin{cases}
 W_i(\operatorname{next}(q,0))\vee
 W_i(\operatorname{next}(q,1)),&\operatorname{owner}(q)=1,\\
 W_i(\operatorname{next}(q,0))\wedge
 W_i(\operatorname{next}(q,1)),&\operatorname{owner}(q)=0.
 \end{cases}
 \tag{S.10}
\]
As in (S.1), primitive recursion on the whole $2^r$-bit table gives
a total numerical definition.  Alternatively, a clocked depth-first
evaluation uses at most $h+1$ frames, each holding an $r$-bit state,
a move index and a return bit, plus the reusable circuit workspace.
Its space is polynomial in the input size because $h$ is unary.
Its time is at most $2^{h+1}$ times a polynomial local cost.

The winning-strategy construction is the one in S10: choose the least
winning child at one's own winning position and retain both children
at the opponent's position; reverse the roles at a losing position.
The entire tree has fewer than $2^{h+1}$ nodes and a primitive-recursive
code bound, for example $2^{(r+4)2^{h+1}}$ with a fixed tree encoding.
Backward induction proves that the chosen strategy wins against every
opposing play.  Bound all states, histories and strategy codes by
these displayed numbers, and the invariant is a bounded arithmetic
formula available to iPRA.

For hardness take a prenex QBF from S12.  A game state stores the
current variable index and the assignment prefix.  The owner is the
quantifier type at that index; the move sets that variable, and the
terminal circuit evaluates the quantifier-free matrix.  Its width,
transition circuits, and horizon are polynomial in the formula length.
Induction over the remaining quantifier prefix identifies game winning
with QBF truth.  This gives the reduction.  The succinct representation
is essential: the theorem concerns a polynomial-size description of
an exponentially large game tree, not an explicitly listed tree.
\end{proof}

\subsection{Formalization obligations particular to this section}

The companion \texttt{lean/SpaceOutline.lean} fixes numerical graph
and QBF data and the total recurrences.  It is an uncompiled outline,
not an asserted completion of these complexity proofs.  The shortest
route to a complete formalization has the following dependencies.

\begin{enumerate}
\item Establish bit-table access, bounded folds, finite path encoding,
      and list deletion at the numerical level.  Prove S01 using a
      bounded path-code invariant, then compile that invariant with
      ProveIt's bounded-formula reflection rather than using an
      induction principle for arbitrary formulas.
\item Define fixed-machine configuration packing, validity, the explicit
      accepting sink, and the local edge predicate.  Prove both trace
      maps in S02, including malformed-code behavior.  The machine code
      and its explicit clock are inputs to a bounded interpreter; an
      unclocked universal evaluator for primitive-recursive syntax is
      never needed.
\item Prove equality of the matrix recurrence (S.4) with the stack
      implementation.  A theorem about a large matrix alone does not
      prove Savitch's workspace bound.  Formalize the live-frame bound
      and the time clock separately.
\item Define the forced-enumeration transcript checker and the canonical
      successful-transcript constructor.  Prove the two implications in
      S07, and carry both universal soundness and constructed completeness
      through the stage induction of S08.  This is where an unjustified
      appeal to ``guessing nonreachability'' must be excluded.
\item Define fresh-variable generation, vector equality and the single-copy
      construction (S.7).  Prove its evaluation recurrence and its syntax
      count independently.  Prove the edge-formula compiler correct before
      combining these facts into S12.
\item Finally construct the actual object-language formulas and proof
      terms.  The target is ProveIt's intuitionistic derivability type
      \texttt{([] : List (Fm PTm 0))} followed by the
      \texttt{PRAAx} turnstile and the coded sentence.  Ambient Lean
      theorems about finite lists or functions are useful semantic
      checks, but do not by themselves inhabit that object-theory type.
\end{enumerate}
